% Code/prompt rendering configuration (% Code/prompt rendering configuration (% Code/prompt rendering configuration (% Code/prompt rendering configuration (\include{lstlisting} from main.tex).
% Two independent mechanisms live here:
%   (1) listings   — ASCII Python/shell code listings in Chapter 4.
%   (2) tcolorbox + fvextra — verbatim Vietnamese system prompts in Appendix B.
%       listings' UTF-8 reader mangles the diacritics (it needed a large
%       literate accent map), so the prompts now use fvextra (native UTF-8)
%       inside a breakable tcolorbox instead — long prompts split across pages.

% --- (1) Python/shell code listings (Chapter 4) ---
\lstset{
  basicstyle=\ttfamily\small,
  breaklines=true,
  frame=single,
  numbers=left,
  numberstyle=\tiny,
  showstringspaces=false,
  tabsize=2,
  captionpos=b,
  language=Python
}
\renewcommand{\lstlistingname}{Listing}

% --- (2) Vietnamese system-prompt boxes (Appendix B) ---
% Symbols that occur in the prompts but are not set up for direct UTF-8 use
% (others, e.g. em dash, are predefined; declared here so verbatim is robust).
\newunicodechar{→}{\ensuremath{\rightarrow}}
\newunicodechar{×}{\ensuremath{\times}}
\newunicodechar{—}{\textemdash}

% Verbatim body: native UTF-8 (keeps the diacritics), wraps long lines anywhere.
\DefineVerbatimEnvironment{promptverbatim}{Verbatim}{%
  fontsize=\footnotesize,
  breaklines=true,
  breakanywhere=true,
}

% Breakable, titled, numbered box for each prompt.
% Usage: \begin{promptbox}[label={lst:...}]{<caption>} ... \end{promptbox}
\newtcolorbox[auto counter]{promptbox}[2][]{%
  breakable, enhanced,
  colback=gray!4,
  colframe=black!72,
  coltitle=white,
  fonttitle=\bfseries\small,
  title={Prompt~\thetcbcounter:~#2},
  left=2mm, right=2mm, top=1mm, bottom=1mm,
  #1,
}
 from main.tex).
% Two independent mechanisms live here:
%   (1) listings   — ASCII Python/shell code listings in Chapter 4.
%   (2) tcolorbox + fvextra — verbatim Vietnamese system prompts in Appendix B.
%       listings' UTF-8 reader mangles the diacritics (it needed a large
%       literate accent map), so the prompts now use fvextra (native UTF-8)
%       inside a breakable tcolorbox instead — long prompts split across pages.

% --- (1) Python/shell code listings (Chapter 4) ---
\lstset{
  basicstyle=\ttfamily\small,
  breaklines=true,
  frame=single,
  numbers=left,
  numberstyle=\tiny,
  showstringspaces=false,
  tabsize=2,
  captionpos=b,
  language=Python
}
\renewcommand{\lstlistingname}{Listing}

% --- (2) Vietnamese system-prompt boxes (Appendix B) ---
% Symbols that occur in the prompts but are not set up for direct UTF-8 use
% (others, e.g. em dash, are predefined; declared here so verbatim is robust).
\newunicodechar{→}{\ensuremath{\rightarrow}}
\newunicodechar{×}{\ensuremath{\times}}
\newunicodechar{—}{\textemdash}

% Verbatim body: native UTF-8 (keeps the diacritics), wraps long lines anywhere.
\DefineVerbatimEnvironment{promptverbatim}{Verbatim}{%
  fontsize=\footnotesize,
  breaklines=true,
  breakanywhere=true,
}

% Breakable, titled, numbered box for each prompt.
% Usage: \begin{promptbox}[label={lst:...}]{<caption>} ... \end{promptbox}
\newtcolorbox[auto counter]{promptbox}[2][]{%
  breakable, enhanced,
  colback=gray!4,
  colframe=black!72,
  coltitle=white,
  fonttitle=\bfseries\small,
  title={Prompt~\thetcbcounter:~#2},
  left=2mm, right=2mm, top=1mm, bottom=1mm,
  #1,
}
 from main.tex).
% Two independent mechanisms live here:
%   (1) listings   — ASCII Python/shell code listings in Chapter 4.
%   (2) tcolorbox + fvextra — verbatim Vietnamese system prompts in Appendix B.
%       listings' UTF-8 reader mangles the diacritics (it needed a large
%       literate accent map), so the prompts now use fvextra (native UTF-8)
%       inside a breakable tcolorbox instead — long prompts split across pages.

% --- (1) Python/shell code listings (Chapter 4) ---
\lstset{
  basicstyle=\ttfamily\small,
  breaklines=true,
  frame=single,
  numbers=left,
  numberstyle=\tiny,
  showstringspaces=false,
  tabsize=2,
  captionpos=b,
  language=Python
}
\renewcommand{\lstlistingname}{Listing}

% --- (2) Vietnamese system-prompt boxes (Appendix B) ---
% Symbols that occur in the prompts but are not set up for direct UTF-8 use
% (others, e.g. em dash, are predefined; declared here so verbatim is robust).
\newunicodechar{→}{\ensuremath{\rightarrow}}
\newunicodechar{×}{\ensuremath{\times}}
\newunicodechar{—}{\textemdash}

% Verbatim body: native UTF-8 (keeps the diacritics), wraps long lines anywhere.
\DefineVerbatimEnvironment{promptverbatim}{Verbatim}{%
  fontsize=\footnotesize,
  breaklines=true,
  breakanywhere=true,
}

% Breakable, titled, numbered box for each prompt.
% Usage: \begin{promptbox}[label={lst:...}]{<caption>} ... \end{promptbox}
\newtcolorbox[auto counter]{promptbox}[2][]{%
  breakable, enhanced,
  colback=gray!4,
  colframe=black!72,
  coltitle=white,
  fonttitle=\bfseries\small,
  title={Prompt~\thetcbcounter:~#2},
  left=2mm, right=2mm, top=1mm, bottom=1mm,
  #1,
}
 from main.tex).
% Two independent mechanisms live here:
%   (1) listings   — ASCII Python/shell code listings in Chapter 4.
%   (2) tcolorbox + fvextra — verbatim Vietnamese system prompts in Appendix B.
%       listings' UTF-8 reader mangles the diacritics (it needed a large
%       literate accent map), so the prompts now use fvextra (native UTF-8)
%       inside a breakable tcolorbox instead — long prompts split across pages.

% --- (1) Python/shell code listings (Chapter 4) ---
\lstset{
  basicstyle=\ttfamily\small,
  breaklines=true,
  frame=single,
  numbers=left,
  numberstyle=\tiny,
  showstringspaces=false,
  tabsize=2,
  captionpos=b,
  language=Python
}
\renewcommand{\lstlistingname}{Listing}

% --- (2) Vietnamese system-prompt boxes (Appendix B) ---
% Symbols that occur in the prompts but are not set up for direct UTF-8 use
% (others, e.g. em dash, are predefined; declared here so verbatim is robust).
\newunicodechar{→}{\ensuremath{\rightarrow}}
\newunicodechar{×}{\ensuremath{\times}}
\newunicodechar{—}{\textemdash}

% Verbatim body: native UTF-8 (keeps the diacritics), wraps long lines anywhere.
\DefineVerbatimEnvironment{promptverbatim}{Verbatim}{%
  fontsize=\footnotesize,
  breaklines=true,
  breakanywhere=true,
}

% Breakable, titled, numbered box for each prompt.
% Usage: \begin{promptbox}[label={lst:...}]{<caption>} ... \end{promptbox}
\newtcolorbox[auto counter]{promptbox}[2][]{%
  breakable, enhanced,
  colback=gray!4,
  colframe=black!72,
  coltitle=white,
  fonttitle=\bfseries\small,
  title={Prompt~\thetcbcounter:~#2},
  left=2mm, right=2mm, top=1mm, bottom=1mm,
  #1,
}
